% ---- begin paper/tables/a1_mh.tex
\begin{table}[htbp]\centering
\caption{Mental-Health Instruments: Items and Scales}
\label{tab:a_mh}
\footnotesize\renewcommand{\arraystretch}{1.0}
\begin{threeparttable}
\begin{tabular}{@{}p{68pt}p{212pt}p{150pt}@{}}
\toprule
 & Question, verbatim & Response scale \\
 & (1) & (2) \\
\midrule
\panel{3}{Panel A. PHQ-4, answered by the adolescent in private (2024 wave)}
\hspace{1em}\texttt{d4\_1} & ?`Te has sentido bajoneado(a), deprimido(a), irritable o desesperanzado(a)? & 0 = Nunca; 1 = Algunos días; 2 = Más de la mitad de los días; 3 = Casi todos los días \\
\addlinespace[1pt]
\hspace{1em}\texttt{d4\_2} & ?`Has sentido poco interés o placer al hacer las cosas? & 0 = Nunca; 1 = Algunos días; 2 = Más de la mitad de los días; 3 = Casi todos los días \\
\addlinespace[1pt]
\hspace{1em}\texttt{d4\_3} & ?`Te has sentido muy nervioso(a), angustiado(a) o con los nervios de punta? & 0 = Nunca; 1 = Algunos días; 2 = Más de la mitad de los días; 3 = Casi todos los días \\
\addlinespace[1pt]
\hspace{1em}\texttt{d4\_4} & ?`No has podido dejar de preocuparte o controlar la preocupación? & 0 = Nunca; 1 = Algunos días; 2 = Más de la mitad de los días; 3 = Casi todos los días \\
\addlinespace
\panel{3}{Panel B. Child Behavior Checklist, answered by the main caregiver about the adolescent}
\hspace{1em}\texttt{cbcl2} & Total problems scale. The survey releases the computed scores, not the copyrighted items. & T score, standardized in the analysis \\
\hspace{1em}\texttt{cbcl2 \_i} & Internalizing problems: the anxious/depressed, withdrawn and somatic-complaints syndrome scales. & T score, standardized \\
\hspace{1em}\texttt{cbcl2 \_e} & Externalizing problems: the rule-breaking and aggressive-behavior syndrome scales. & T score, standardized \\
\hspace{1em}\texttt{cbcl2 \_6} & Attention problems syndrome scale. & T score, standardized \\
\bottomrule
\end{tabular}
\begin{wpnotes}
\textit{Notes:} Question wording and response labels are reproduced verbatim in the original Spanish of the questionnaire, from the variable and value labels of the public microdata; wording the microdata truncates at 80 characters is marked with an ellipsis. The PHQ-4 score used in the paper is the sum of the four items rescaled to 0--3 each (0--12 in total) and standardized within age in years. The positive PHQ-2 (depression) and GAD-2 (anxiety) screens are the survey-provided indicators for a subscore of three or more on the corresponding pair of items.
\end{wpnotes}
\end{threeparttable}
\end{table}
% ---- end paper/tables/a1_mh.tex
